\documentclass[../../main.tex]{subfiles}
\begin{document}

{\bf [解]}

% TODO: upgrade
% TODO: upgrade
\begin{figure}[htb]
  \centering
  % =========================================================================
  % 左図: 正三角形の3分割領域 P, Q, R
  % =========================================================================
  \begin{tikzpicture}[scale=1.5, >=stealth, every node/.style={font=\small}]

    % 正三角形の頂点 (一辺の長さを 2*sqrt(3), 外接円半径 2)
    % A: 上頂点 (0, 2), B: 左下 (-sqrt(3), -1), C: 右下 (sqrt(3), -1)
    \coordinate (A) at (0, 2);
    \coordinate (B) at (-1.732, -1);
    \coordinate (C) at (1.732, -1);
    \coordinate (G) at (0, 0); % 重心

    % 各辺の中点
    \coordinate (M_BC) at (0, -1);
    \coordinate (M_CA) at (0.866, 0.5);
    \coordinate (M_AB) at (-0.866, 0.5);

    % ------------------------------------------------------------
    % 3つの領域 P, Q, R の塗り分け
    % P: 上側頂点 A を含む凧形 (A - M_CA - G - M_AB)
    % Q: 右下頂点 C を含む凧形 (C - M_BC - G - M_CA)
    % R: 左下頂点 B を含む凧形 (B - M_AB - G - M_BC)
    % ------------------------------------------------------------
    \fill[blue!15, opacity=0.7] (A) -- (M_CA) -- (G) -- (M_AB) -- cycle;
    \fill[red!15, opacity=0.7]  (C) -- (M_BC) -- (G) -- (M_CA) -- cycle;
    \fill[teal!15, opacity=0.7] (B) -- (M_AB) -- (G) -- (M_BC) -- cycle;

    % 正三角形の外枠
    \draw[line width=1.4pt, gray!20!black] (A) -- (B) -- (C) -- cycle;

    % 中線・境界線 (各頂点から対辺の中点へ)
    \draw[thick, gray!50!black] (A) -- (M_BC);
    \draw[thick, gray!50!black] (B) -- (M_CA);
    \draw[thick, gray!50!black] (C) -- (M_AB);

    % 重心 G
    \fill[black] (G) circle (1.2pt);

    % 領域ラベル
    \node[blue!90!black, font=\large\bfseries] at (0, 0.75) {$P$};
    \node[red!90!black, font=\large\bfseries]  at (0.6, -0.45) {$Q$};
    \node[teal!90!black, font=\large\bfseries] at (-0.6, -0.45) {$R$};

    % 図タイトル
    \node[below=14pt, font=\footnotesize\bfseries] at (0, -1.0) {(a) 正三角形の対称な 3 領域};

  \end{tikzpicture}
  \qquad\quad
  % =========================================================================
  % 右図: 2ステップごとの状態遷移図 (P, Q, R)
  % =========================================================================
  \begin{tikzpicture}[scale=1.4, >=stealth, every node/.style={font=\small}]

    % 状態ノード (円形ノードで配置)
    \node[draw=blue!85!black, fill=blue!8, circle, line width=1.2pt, inner sep=4pt, font=\normalsize\bfseries]
    (P) at (0, 1.4) {$P$};
    \node[draw=red!85!black, fill=red!8, circle, line width=1.2pt, inner sep=4pt, font=\normalsize\bfseries]
    (Q) at (1.8, -0.8) {$Q$};
    \node[draw=teal!85!black, fill=teal!8, circle, line width=1.2pt, inner sep=4pt, font=\normalsize\bfseries]
    (R) at (-1.8, -0.8) {$R$};

    % 遷移矢印 (Q への流入を中心に描画)
    % P -> Q (確率 1/6)
    \draw[->, line width=1.2pt, blue!85!black] (P) -- (Q)
    node[midway, above right=1pt, font=\footnotesize\bfseries] {$\frac{1}{6}$};

    % R -> Q (確率 1/6)
    \draw[->, line width=1.2pt, teal!85!black] (R) -- (Q)
    node[midway, below=2pt, font=\footnotesize\bfseries] {$\frac{1}{6}$};

    % Q -> Q (自己ループ, 確率 2/3)
    \draw[->, line width=1.2pt, red!85!black] (Q) to[out=-25, in=45, looseness=4.5]
    node[right=2pt, font=\footnotesize\bfseries] {$\frac{2}{3}$} (Q);

    % 対称性の補助（破線で他の同等な遷移を示唆）
    \draw[->, dashed, line width=0.9pt, gray!50] (P) -- (R)
    node[midway, above left=1pt, font=\scriptsize] {$\frac{1}{6}$};
    \draw[->, dashed, line width=0.9pt, gray!50] (R) to[out=135, in=-155, looseness=4.5]
    node[left=2pt, font=\scriptsize] {$\frac{2}{3}$} (R);
    \draw[->, dashed, line width=0.9pt, gray!50] (P) to[out=120, in=60, looseness=4.5]
    node[above=2pt, font=\scriptsize] {$\frac{2}{3}$} (P);

    % 図タイトル
    \node[below=14pt, font=\footnotesize\bfseries] at (0, -1.0) {(b) $2$ 秒ごとの状態遷移確率};

  \end{tikzpicture}
  \caption{正三角形内の対称な領域分割と球の状態遷移図}
  \label{fig:triangle_markov_transition_upgrade}
\end{figure}

上図のように$R$ を定める．
$k \in \mathbb{N}$ とすると球は $2k$ 秒後にのみ $P, Q, R$ にある．
対称性から，$2k$ 秒後に球が $P, Q, R$ にある確率は $a_k, b_k, b_k$ とおける．
題意から以下の漸化式をうる．
\begin{align}
  b_{k+1} = \dfrac{2}{3} b_k + \dfrac{1}{6} b_k + \dfrac{1}{6} a_k \label{eq:1}
\end{align}
ここで，$2k$ 秒後には球は$P, Q, R$のいずれかにあるから
\begin{align}
  a_k + 2b_k = 1 \label{eq:2}
\end{align}
である．\cref{eq:1}に代入して$a_k$ を消去すると
\begin{align}
  & b_{k+1} = \dfrac{1}{2} b_k + \dfrac{1}{6} \\
  & b_{k+1} - \dfrac{1}{3} = \dfrac{1}{2} \left( b_k - \dfrac{1}{3} \right)
\end{align}
である．初期条件は$b_0 = 0$ だから，漸化式を繰り返し用いて
\begin{align}
  b_k = -\dfrac{1}{3}\left(\dfrac{1}{2}\right)^k + \dfrac{1}{3}
\end{align}
だから，もとめる確率は
\begin{align}
  \begin{cases}
    0 & (n \in \text{odd}) \\
    \dfrac{1}{3} \left( 1 - \left(\dfrac{1}{2}\right)^{\frac{n}{2}} \right) & (n \in \text{even})
  \end{cases}
\end{align}
である．

\bigskip

\end{document}
